\chapter*{Conclusiones}
\addcontentsline{toc}{chapter}{Conclusiones}\label{chap:conclusions}

El trabajo abordó el desarrollo de un sistema OCR para texto impreso en español y su integración en una aplicación web de digitalización documental, en respuesta a la necesidad del Instituto Cubano de Investigación Cultural Juan Marinello de procesar el acervo documental que la institución conserva sobre soporte físico. La solución articula cuatro componentes técnicos —generador de \textit{corpus} sintético, \textit{pipeline} de preprocesamiento, modelo de reconocimiento de secuencias y aplicación web— cuyos resultados resume el listado siguiente.

\begin{enumerate}[leftmargin=*,labelsep=0.5em]
	\item \textbf{Análisis del estado del arte y decisiones de diseño.} El estudio comparativo entre motores OCR clásicos, arquitecturas que operan sobre aprendizaje profundo y plataformas de digitalización documental estableció que la combinación de una red neuronal convolucional recurrente con la función de pérdida CTC ofrece el equilibrio adecuado entre coste de entrenamiento y capacidad de modelado del contexto léxico para el reconocimiento de líneas de texto impreso. La revisión de los conjuntos de datos públicos para OCR en español identificó cuatro limitaciones que descartaban su uso como base de entrenamiento —cobertura incompleta de los caracteres específicos del español, redundancia tipográfica, inviabilidad de corrección por augmentación posterior y ausencia de contexto léxico—, y justificó la generación de un \textit{corpus} sintético propio.
	
	\item \textbf{Corpus sintético.} El trabajo diseñó e implementó un generador de pares imagen--transcripción con control explícito sobre la distribución de frecuencias de caracteres mediante una estrategia de sobremuestreo en tres grupos. El generador toma el texto de obras literarias en español del dominio público disponibles en Project Gutenberg, lo renderiza con un repertorio de quince fuentes tipográficas que cubren tres familias estilísticas (monoespacio, serif clásica y máquina de escribir), y aplica una simulación de degradación visual compatible con los dos métodos de binarización del \textit{pipeline}. El \textit{corpus} final reúne 19\,995 imágenes con cobertura total del vocabulario de cien símbolos, y mantiene todos los caracteres específicos del español formal —vocales con tilde, eñe, comillas angulares, raya de diálogo y signos de apertura de interrogación y exclamación— por encima del umbral mínimo de quinientas apariciones.
	
	\item \textbf{Pipeline de preprocesamiento.} El trabajo implementó un \textit{pipeline} con configuración automática a partir de cinco métricas de imagen, sin intervención del usuario. El \textit{pipeline} elige entre la binarización global de Otsu y la binarización adaptativa de Sauvola según la uniformidad del fondo, corrige la inclinación en tres niveles (global de la página, residual por bloque y curvatura por línea), detecta bloques de texto mediante una variante recursiva del algoritmo XY-Cut robusta a páginas con doble columna y a pliegos abiertos, y segmenta las líneas mediante proyección horizontal de la imagen binaria. La salida del \textit{pipeline} son líneas con altura fija de sesenta y cuatro píxeles, en el formato que el modelo de reconocimiento exige.
	
	\item \textbf{Modelo de reconocimiento y protocolo de evaluación.} El trabajo diseñó, entrenó y evaluó un modelo CRNN con red recurrente bidireccional y función de pérdida CTC. La comparación durante el entrenamiento entre decodificación \textit{greedy} y búsqueda en haz pura mostró que la \textit{greedy} alcanza un CER inferior sobre la validación sintética, dada la alta confianza del modelo visual sobre imágenes tipográficamente limpias. La integración posterior de un modelo de lenguaje de cinco gramas a nivel de carácter, que el equipo entrenó previamente sobre \textit{corpus} de español de uso general y que la herramienta KenLM carga, dota al \textit{decoder} de un criterio adicional para reordenar hipótesis acústicamente compatibles. El sistema en producción combina la búsqueda en haz con este modelo de lenguaje, configuración que mejora el rendimiento sobre documentos reales gracias a la mayor incertidumbre acústica del modelo visual ante degradaciones de escaneado y a la regularidad morfológica del español. La inspección cualitativa sobre líneas reales del archivo del centro Juan Marinello, que la subsección~\ref{subsec:lm-posterior} documenta, confirma este principio y caracteriza el tipo de correcciones que el LM introduce: eliminación de sufijos léxicos espurios, limpieza de puntuación ortográficamente incorrecta y sustitución de no-palabras por palabras del idioma. El protocolo de evaluación empleó las métricas estándar del campo —tasa de error de carácter, tasa de error de palabra y exactitud por línea—, intervalos de confianza al~95\,\% que el \textit{bootstrap} produce sobre diez mil remuestreos del conjunto de validación, pruebas no paramétricas para comparación entre configuraciones y experimentos de exclusión de una fuente, que la literatura inglesa denomina \textit{Leave-One-Font-Out}, para cuantificar la generalización a tipografías que quedan fuera del entrenamiento. Tras la etapa de \textit{fine-tuning} sobre el \textit{corpus} extendido de quince fuentes, el modelo final alcanzó una tasa de error de carácter del~0,88\,\%, una tasa de error de palabra del~3,31\,\% y una exactitud por línea del~83,04\,\% sobre el conjunto de validación, con uniformidad de rendimiento entre las quince fuentes que la prueba de Kruskal--Wallis confirma ($H = 17{,}04$, $p = 0{,}254$). El trabajo realizó la evaluación final del modelo desplegado sobre un conjunto de prueba independiente de 1\,995 imágenes que el mismo procedimiento de generación produce con una semilla distinta y que el protocolo reserva en exclusiva para una única medición al término del \textit{fine-tuning}; el modelo alcanzó sobre este conjunto una tasa de error de carácter del~4,28\,\% (IC 95\,\% de $[0{,}0407,\ 0{,}0448]$), una tasa de error de palabra del~9,73\,\% y una exactitud por línea del~37,4\,\%, valores que constituyen la estimación insesgada del rendimiento que cabe esperar del modelo en condiciones operativas sobre material sintético comparable al del entrenamiento.
	
	\item \textbf{Aplicación web.} El trabajo desarrolló una aplicación web en Django que integra el sistema OCR y ofrece un visor de doble panel con el facsimilar original y la transcripción, junto con funcionalidades de búsqueda, gestión documental y revisión colaborativa con roles diferenciados. El sistema permite exportar las transcripciones en los formatos estándar de la comunidad de digitalización documental: ALTO~\cite{loc-alto}, PAGE-XML~\cite{pletschacher2010page} y TEI~\cite{tei-consortium}.
	
	\item \textbf{Síntesis.} La combinación de los tres elementos distintivos del trabajo —un \textit{corpus} sintético anclado en las frecuencias estadísticas del español, un modelo CRNN+CTC que aprende sobre ese \textit{corpus}, y una aplicación web con revisión colaborativa— constituye un sistema funcional para la digitalización del acervo del Instituto Cubano de Investigación Cultural Juan Marinello. El trabajo cumple en su totalidad los seis objetivos específicos que la introducción enuncia.
\end{enumerate}

\noindent\textbf{Limitaciones del trabajo.} El sistema atiende texto impreso: manuscritos, documentos mecanografiados con desgaste fuera del rango que la simulación de degradación cubre, y tipografías con perfiles visuales distantes de las quince fuentes del \textit{corpus} quedan fuera del alcance del trabajo. Los documentos con diseño visual complejo —revistas en las que los bloques de texto enlazan unos con otros, anotaciones a mano sobre impreso, marginalia, tablas— requieren una etapa adicional de detección del orden de lectura que el \textit{pipeline} actual no incorpora. La inferencia opera sobre líneas individuales, y el sistema no aprovecha el contexto entre líneas consecutivas, lo que podría aportar precisión adicional en pasajes con vocabulario poco frecuente. El modelo de lenguaje de cinco gramas que opera en producción captura dependencias locales —sufijos, terminaciones verbales, reglas ortográficas obligatorias del español— pero no captura dependencias a larga distancia que un modelo de lenguaje neuronal sí capturaría. La evaluación reportada sobre el conjunto de prueba sintético constituye una estimación insesgada del rendimiento del modelo en condiciones controladas, sobre material léxico independiente del entrenamiento pero generado por el mismo procedimiento; el rendimiento sobre documentos reales digitalizados, con la heterogeneidad de calidad de escaneo, los artefactos físicos del soporte y la variabilidad tipográfica del archivo histórico que estos presentan, recibe en este trabajo una caracterización cualitativa sobre líneas del archivo del centro Juan Marinello pero queda pendiente de cuantificación rigurosa sobre un \textit{benchmark} con transcripciones de referencia validadas humanamente, lo que constituye una de las líneas naturales de extensión del trabajo.

\noindent\textbf{Dificultades encontradas.} La principal dificultad del trabajo fue la escasez de conjuntos de datos públicos para OCR en español que cubrieran los signos específicos del idioma y que tuvieran un tamaño comparable al de los recursos en inglés. La decisión de generar un \textit{corpus} sintético propio resolvió la dificultad pero exigió diseñar e implementar un generador con control fino sobre la distribución de frecuencias de caracteres, la selección tipográfica y la simulación de degradación visual. Una dificultad secundaria fue la calibración de los rangos de ruido gaussiano, ruido sal y pimienta y suavizado, que el desarrollo ajustó de forma iterativa para preservar la compatibilidad con los dos métodos de binarización del \textit{pipeline} sin perder caracteres ni introducir artefactos que el modelo pudiera memorizar como rasgo discriminativo.
